\Needspace{15\baselineskip}
\fancyhead[L]{\small JUNO ELEMENTS VIII.}
\begin{center}
JUNO ELEMENTS VIII.
\end{center}

\hrule
\begin{center}
\fontsize{10}{12}\selectfont Göttingische gelehrte Anzeigen. Issue 136. Pp. 1353..1355. 25 August 1808.
\end{center}
\par\vspace{4pt}\hrule\vspace{4pt}

\par\indent The first observations of \emph{Juno} made at our observatory this year have already been reported in issue 107 of this journal [p. 303 of this volume]. Observations have continued whenever conditions were favorable, though only with the ring micrometer, since the small planet remained invisible throughout in our mural quadrant’s telescope, which gathers little light. An effort was made to compensate for this limitation with many comparisons, carried out as carefully as possible. The four observations near opposition in particular may be considered as accurate as the method allows. Since the observation of 22 June has undergone a small correction, we give the first two determinations again here along with all the later ones made so far [they appear above on page 305 in the letter to \textsc{von Lindenau}, 1808 Aug. 6].

\par\indent The latest observations thus show that the error in the ephemeris has grown to 13 minutes in right ascension. The error in declination has hardly changed and is still about $1\frac{1}{2}$ minutes. No observations from elsewhere have become known so far.

\par\indent Professor \textsc{Gauss} used the last four observations, which turned out particularly well, to determine the opposition. He obtained the following result:

\begin{center}
1808. August 2. $9^{\text{u}}\,30^{\text{m}}\,43^{\text{s}}$ mean time at Göttingen\\
true longitude .\;.\;.\;.\;.\;.\;.\;.\;.\; $310^\circ\,16'\,31''5$\\
true geocentric latitude .\;.\;.\;.\; $13\;53\;56,2$ north
\end{center}

\Needspace{25\baselineskip}
\par\indent This opposition, the \emph{fourth} observed so far, was then combined with those of 1804, 1806, and 1807 to determine the following new elements (VIII), which still fit all observations made to date very well.


\begin{center}\fontsize{10}{12}\selectfont\setlength{\tabcolsep}{3pt}\renewcommand{\arraystretch}{1.15}\begin{tabular}{lr}
\multicolumn{2}{c}{Epoch of mean longitude}\\
\multicolumn{2}{c}{for the meridian of Göttingen:}\\\hline
1804 .\;.\;.\;.\;.\;.\;.\;.\; & $320^\circ\;1'\;20''1$\\
1805 .\;.\;.\;.\;.\;.\;.\;.\; & $42\;35\;8.4$\\
1806 .\;.\;.\;.\;.\;.\;.\;.\; & $125\;8\;56.7$\\
1807 .\;.\;.\;.\;.\;.\;.\;.\; & $207\;42\;45.0$\\
1808 .\;.\;.\;.\;.\;.\;.\;.\; & $290\;30\;7.6$\\
1809 .\;.\;.\;.\;.\;.\;.\;.\; & $13\;3\;55.8$\\
1810 .\;.\;.\;.\;.\;.\;.\;.\; & $95\;37\;44.1$\\
\end{tabular}\end{center}

\begin{center}\fontsize{9}{11}\selectfont\setlength{\tabcolsep}{3pt}\renewcommand{\arraystretch}{1.15}\begin{tabular}{lr}
Daily mean tropical motion .\;.\; & $814''324$\\
Tropical orbital period .\;.\; & 1591 days 12 hours\\
Perihelion 1805 .\;.\;.\;.\;.\;.\;.\; & $53^\circ\;10'\;53''9$\\
Ascending node 1805 .\;.\;.\;.\; & $171\;4\;11{,}3$\\
\multicolumn{2}{l}{\quad Both assumed fixed with respect to the stars.}\\
Inclination of the orbit .\;.\;.\;.\;.\;.\;.\;.\; & $13\;4\;11{,}0$\\
Eccentricity .\;.\;.\;.\;.\;.\;.\;.\;.\;.\;.\; & $0{,}2554521$\\
Logarithm of the semimajor axis .\;.\;.\; & $0{,}4261883$\\
\end{tabular}\end{center}


\SourceLeaf{312}\par\indent According to these elements, the next opposition, when Juno will be considerably brighter again, falls on 30 January 1810, at longitude $130^\circ\,2'$ and south latitude $14^\circ\,51'$, at the head of Hydra.

